\ifdefined\gkver\else\def\gkver{student}\fi
\documentclass[\gkver]{gaokaozhenti}
\begin{document}

\chapter{2003年全国春季理综}
%% paperType: 理综物理部分

\begin{enumerate}
\item
%% number: 15
%% paperName: 2003年全国春季理综第 15 题 6 分
%% typeId: 1
%% type: 单选题
%% chapter: 机械振动
%% point: 振动图象
%% method: 无
%% score: 6
%% degree: 950
%% duplicateId: 0
%% body:
右图表示一简谐横波波源的振动图象。根据图象可确定该波的\xzanswer{D}

\onepicture[w=7.2cm]{figs/15.png}

\fourchoices[answer=D]
{波长，波速}
{周期，波速}
{波长，振幅}
{周期，振幅}

%% answer: D
%% memo:
\memoanswer{
本题原卷及材料中未提供详解，待补充。
}

\item
%% number: 16
%% paperName: 2003年全国春季理综第 16 题 6 分
%% typeId: 1
%% type: 单选题
%% chapter: 运动和力
%% point: 牛顿第三定律
%% method: 无
%% score: 6
%% degree: 550
%% duplicateId: 0
%% body:
在滑冰场上，甲、乙两小孩分别坐在滑冰板上，原来静止不动，在相互猛推一下后分别向相反方向运动。假定两板与冰面间的摩擦因数相同。已知甲在冰上滑行的距离比乙远，这是由于\xzanswer{C}

\fourchoices[answer=C]
{在推的过程中，甲推乙的力小于乙推甲的力}
{在推的过程中，甲推乙的时间小于乙推甲的时间}
{在刚分开时，甲的初速度大于乙的初速度}
{在分开后，甲的加速度的大小小于乙的加速度的大小}

%% answer: C
%% memo:
\memoanswer{
本题原卷及材料中未提供详解，待补充。
}

\item
%% number: 17
%% paperName: 2003年全国春季理综第 17 题 6 分
%% typeId: 1
%% type: 单选题
%% chapter: 原子和原子核
%% point: 质能方程
%% method: 无
%% score: 6
%% degree: 750
%% duplicateId: 0
%% body:
下面是一核反应方程 \ce{^2_1 H + ^3_1 H -> ^4_2 He + X}，用 $c$ 表示光速，则\xzanswer{D}

\fourchoices[answer=D]
{X 是质子，核反应放出的能量等于质子质量乘 $c^{2}$}
{X 是中子，核反应放出的能量等于质子质量乘 $c^{2}$}
{X 是质子，核反应放出的能量等于氘核与氚核的质量和减去氦核与质子的质量和，再乘 $c^{2}$}
{X 是中子，核反应放出的能量等于氘核与氚核的质量和减去氦核与中子的质量和，再乘 $c^{2}$}

%% answer: D
%% memo:
\memoanswer{
本题原卷及材料中未提供详解，待补充。
}

\item
%% number: 18
%% paperName: 2003年全国春季理综第 18 题 6 分
%% typeId: 1
%% type: 单选题
%% chapter: 气体性质
%% point: 玻意耳定律
%% method: 无
%% score: 6
%% degree: 650
%% duplicateId: 0
%% body:
一根粗细均匀长 $1.0\Um$ 的直玻璃管，上端封闭，下端开口，将它竖直地缓慢插入深水池中，直到管内水面距管上端 $0.50\Um$ 为止。已知水的密度为 $1.0\times10^{3}\Ukgmc$，重力加速度为 $10\Umsq$，$1.0\times10^{5}\UPa$，则这时管内、外水面的高度差为\xzanswer{C}

\fourchoices[answer=C]
{$9\Um$}
{$9.5\Um$}
{$10\Um$}
{$10.5\Um$}

%% answer: C
%% memo:
\memoanswer{
本题原卷及材料中未提供详解，待补充。
}

\item
%% number: 19
%% paperName: 2003年全国春季理综第 19 题 6 分
%% typeId: 1
%% type: 单选题
%% chapter: 光学
%% point: 光的折射
%% method: 无
%% score: 6
%% degree: 650
%% duplicateId: 0
%% body:
如图所示，一玻璃棱镜的横截面是等腰 $\triangle abc$，其中 ac 面是镀银的。现有一光线垂直于 ab 面入射，在棱镜内经过两次反射后垂直于 bc 面射出。则\xzanswer{D}

\onepicture[h=5.5cm]{figs/19.png}

\fourchoices[answer=D]
{$\angle a=30^{\circ}$，$\angle b=75^{\circ}$}
{$\angle a=32^{\circ}$，$\angle b=74^{\circ}$}
{$\angle a=34^{\circ}$，$\angle b=73^{\circ}$}
{$\angle a=36^{\circ}$，$\angle b=72^{\circ}$}

%% answer: D
%% memo:
\memoanswer{
本题原卷及材料中未提供详解，待补充。
}

\item
%% number: 20
%% paperName: 2003年全国春季理综第 20 题 6 分
%% typeId: 1
%% type: 单选题
%% chapter: 曲线运动
%% point: 人造地球卫星
%% method: 无
%% score: 6
%% degree: 650
%% duplicateId: 0
%% body:
在地球（看作质量均匀分布的球体）上空有许多同步卫星，下面说法中正确的是\xzanswer{A}

\fourchoices[answer=A]
{它们的质量可能不同}
{它们的速度可能不同}
{它们的向心加速度可能不同}
{它们离地心的距离可能不同}

%% answer: A
%% memo:
\memoanswer{
本题原卷及材料中未提供详解，待补充。
}

\item
%% number: 27
%% paperName: 2003年全国春季理综第 27 题 18 分
%% typeId: 4
%% type: 实验题
%% chapter: 电路
%% point: 多用表的使用
%% method: 无
%% score: 18
%% degree: 550
%% duplicateId: 0
%% body:
根据试题要求填空或作图。

\twopicture[numA=1, numB=2, labelA=2003全国春季理综27a, labelB=2003全国春季理综27b, wA=3.4cm, wB=4.2cm]
{figs/27a.png}
{figs/27b.png}

\begin{enumerate}
	\item
	图 1 为多用电表的示意图，其中 S、K、T 为三个可调节的部件，现用此电表测量一阻值约为 $20\sim30\UO$ 的定值电阻，测量的某些操作步骤如下：
	\begin{steps}
		\item
		调节可调部件 \tkanswer[2.0cm]{S}，使电表指针停在 \tkanswer[3.0cm]{左端的“0”} 位置；
		\item
		调节可调部件 K，使它的尖端指向 \tkanswer[2.6cm]{$\UO$ 挡的 $\times1$} 位置；
		\item
		将红、黑表笔分别插入“$+$”“$-$”插孔，笔尖相互接触，调节可调部件 \tkanswer[2.0cm]{T}，使电表指针指向 \tkanswer[3.0cm]{右端的“0”} 位置。
	\end{steps}
	\item
	在用多用表测量另一电阻的阻值时，电表的读数如图 2 所示，该电阻的阻值为 \tkanswer[2.0cm]{20.0k}$\UO$。
	\item
	现用伏安法测量（1）问中的那个定值电阻 $R_{x}$ 的阻值（约为 $20\sim30\UO$）。所供器材如下：

	电压表 V（量程 $0\sim15\UV$，内阻 $20\UO$），电流表 A$_{1}$（量程 $0\sim50\UmA$，内阻 $20\UO$），电流表 A$_{2}$（量程 $0\sim30\UmA$，内阻 $4\UO$），滑动变阻器 $R$（最大阻值为 $50\UO$，额定电流为 $1\UA$），直流电源 $E$（电动势约 $9\UV$，内阻约 $0.5\UO$），电键 K 连线用的导线若干根。在以上器材中选出适当的器材，画出电路图，要求在图上标出所选器材的符号。
	\item
	该电阻是由电阻丝绕成的，为了求得该电阻线材料的电阻率，需用螺旋测微器测量电阻丝的直径，结果如图 3，其读数为 \tkanswer[2.6cm]{$0.700\Umm$}。

	\onepicture[num=3, label=2003全国春季理综27c, h=3.0cm, align=right]{figs/27c.png}
\end{enumerate}

%% answer: 
\jdanswer{
\begin{enumerate}
	\item
	①S，左端的“0”；②$\UO$ 挡的 $\times1$；③T，右端的“0”。
	\item
	20.0k（写成 20k 给分）。
	\item
	电路图如图所示。

	\onepicture[w=5.0cm]{figs/27_an.png}
	\item
	$0.700\Umm$。
\end{enumerate}
}

%% memo:
\memoanswer{
本题原卷及材料中未提供详解，待补充。
}

\item
%% number: 28
%% paperName: 2003年全国春季理综第 28 题 24 分
%% typeId: 6
%% type: 计算题
%% chapter: 曲线运动
%% point: 平抛运动
%% method: 无
%% score: 24
%% degree: 550
%% duplicateId: 0
%% body:
有一炮竖直向上发射炮弹，炮弹的质量为 $M=6.0\Ukg$（内含炸药的质量可以忽略不计），射出的初速度 $v_{0}=60\Ums$。当炮弹到达最高点时爆炸为沿水平方向运动的两片，其中一片质量为 $m=4.0\Ukg$。现要求这一片不能落到以发射点为圆心、以 $R=600\Um$ 为半径的圆周范围内，则刚爆炸完时两弹片的总动能至少多大？（$g=10\Umsq$，忽略空气阻力）

%% answer: 
\jdanswer{
$E_{k}=6.0\times10^{4}\UJ$
}

%% memo:
\memoanswer{
设炮弹上升到达最高点的高度为 $H$，根据匀变速直线运动规律，有

$v_{0}^{2}=2gH$

设质量为 $m$ 的弹片刚爆炸后的速度为 $V$，另一块的速度为 $v$，根据动量守恒定律，有

$mV=(M-m)v$

设质量为 $m$ 的弹片运动的时间为 $t$，根据平抛运动规律，有

$H=\frac{1}{2}gt^{2}$

$R=Vt$

炮弹刚爆炸后，两弹片的总动能

$E_{k}=\frac{1}{2}mV^{2}+\frac{1}{2}(M-m)v^{2}$

$E_{k}=\frac{1}{2}\frac{MmR^{2}g^{2}}{(M-m)v_{0}^{2}}$

解以上各式，代入数值得

$E_{k}=6.0\times10^{4}\UJ$
}

\item
%% number: 29
%% paperName: 2003年全国春季理综第 29 题 28 分
%% typeId: 6
%% type: 计算题
%% chapter: 磁场
%% point: 带电粒子在复合场中的运动
%% method: 分情况讨论
%% score: 28
%% degree: 550
%% duplicateId: 0
%% body:
如图所示，在 $Oxyz$ 坐标系所在的空间中，可能存在匀强电场或磁场，也可能两者都存在或都不存在。但如果两者都存在，已知磁场平行于 $xy$ 平面。现有一质量为 $m$ 带正电 $q$ 的点电荷沿 $z$ 轴正方向射入此空间中，发现它做速度为 $v_{0}$ 的匀速直线运动。若不计重力，试写出电场和磁场的分布有哪几种可能性。要求对每一种可能性，都要说出其中能存在的关系。不要求推导或说明理由。

\onepicture[w=6.0cm, align=right]{figs/29.png}

%% answer: 
\jdanswer{
以 $E$ 和 $B$ 分别表示电场强度和磁感强度，有以下几种可能：
\begin{enumerate}
	\item
	$E=0$，$B=0$。
	\item
	$E=0$，$B\neq0$。$B$ 的方向与 $z$ 轴的正方向平行或反平行，$B$ 的大小可任意。
	\item
	$E\neq0$，$B\neq0$。磁场方向可在平行于 $xy$ 平面的任何方向，电场 $E$ 方向平行于 $xy$ 平面，并与 $B$ 的方向垂直；当迎着 $z$ 轴正方向看时，由 $B$ 的方向沿顺时针转 $90^{\circ}$ 后就是 $E$ 的方向；$E$ 和 $B$ 的大小可取满足关系式 $\frac{E}{B}=v_{0}$ 的任何值。
\end{enumerate}
}

%% memo:
\memoanswer{
以 $E$ 和 $B$ 分别表示电场强度和磁感强度，有以下几种可能：

（1）$E=0$，$B=0$；

（2）$E=0$，$B\neq0$。$B$ 的方向与 $z$ 轴的正方向平行或反平行。$B$ 的大小可任意；

（3）$E\neq0$，$B\neq0$。磁场方向可在平行于 $xy$ 平面的任何方向；电场 $E$ 方向平行于 $xy$ 平面，并与 $B$ 的方向垂直；当迎着 $z$ 轴正方向看时，由 $B$ 的方向沿顺时针转 $90^{\circ}$ 后就是 $E$ 的方向；$E$ 和 $B$ 的大小可取满足关系式 $\frac{E}{B}=v_{0}$ 的任何值。
}

\item
%% number: 30
%% paperName: 2003年全国春季理综第 30 题 14 分
%% typeId: 6
%% type: 计算题
%% chapter: 电磁感应
%% point: 切割磁感线电动势
%% method: 无
%% score: 14
%% degree: 550
%% duplicateId: 0
%% body:
图 1 是一台发电机定子中的磁场分布图，其中 N、S 是永久磁铁的两个磁极，它们的表面呈半圆柱面形状。M 是圆柱形铁芯，它与磁极的柱面共轴。磁极与铁芯之间的缝隙中形成方向沿圆柱半径、大小近似均匀的磁场，磁感强度 $B=0.050\UT$。

图 2 是该发电机转子的示意图（虚线表示定子的铁芯 M）。矩形线框 abcd 可绕过 ad、cb 边的中点并与图 1 中的铁芯 M 共轴的固定转轴 oo$'$ 旋转，在旋转过程中，线框的 ab、cd 边始终处在图 1 所示的缝隙内的磁场中。已知 ab 边长 $l_{1}=25.0\Ucm$，ad 边长 $l_{2}=10.0\Ucm$，线框共有 $N=8$ 匝导线，放置的角速度 $\omega=250\Urads$。将发电机的输出端接入图中的装置 K 后，装置 K 能使交流电变成直流电，而不改变其电压的大小。直流电的另一个输出端与一可变电阻 $R$ 相连，可变电阻的另一端 P 是直流电的正极，直流电的另一个输出端 Q 是它的负极。

图 3 是可用于测量阿伏加德罗常数的装置示意图，其中 A、B 是两块纯铜片，插在 \ce{CuSO4} 稀溶液中，铜片与引出导线相连，引出端分别为 x、y。

Ⅰ.（14 分）现把直流电的正、负极与两铜片的引线端相连，调节 $R$，使 \ce{CuSO4} 溶液中产生 $I=0.21\UA$ 的电流。假设发电机的内阻可忽略不计，两铜片间的电阻 $r$ 是恒定的。
\begin{enumerate}
	\item
	求每匝线圈中的感应电动势的大小。
	\item
	求可变电阻 $R$ 与 A、B 间电阻 $r$ 之和。
\end{enumerate}

Ⅱ.（9 分）
\begin{enumerate}
	\item
	当以 $I=0.21\UA$ 的电流电解 60 分钟后，测得铜片 A 的质量增加了 $0.25\Ug$，则图 3 装置的 x 端应与直流电的 \tkanswer[1.5cm]{负} 极相连，它是电解池的 \tkanswer[1.5cm]{阴} 极。
	\item
	电解后铜片 B 的质量 \tkanswer[2.0cm]{减少}。（答增加、减少或不变）
	\item
	列式计算实验测得的阿伏加德罗常数 $N_{A}$。（已知电子电量 $e=1.60\times10^{-19}\UC$）
\end{enumerate}

\threepicture[numA=1, numB=2, numC=3, labelA=2003全国春季理综30a, labelB=2003全国春季理综30b, labelC=2003全国春季理综30c, h=4.0cm, align=right]
{figs/30a.png}
{figs/30b.png}
{figs/30c.png}

%% answer: 
\jdanswer{
\begin{enumerate}
	\item
	$\varepsilon=0.31\UV$。
	\item
	$R+r=12\UO$。
\end{enumerate}
\begin{enumerate}
	\item
	负，阴。
	\item
	减少。
	\item
	$N_{A}=\frac{64\Ugmol\times0.21\UA\times3600\Us}{0.25\Ug\times2\times1.6\times10^{-19}\UC}=6.0\times10^{23}\Umoln$。
\end{enumerate}
}

%% memo:
\memoanswer{
Ⅰ

（1）设线框边的速度为 $v$，则

$v=\frac{1}{2}l_{2}\omega$

一匝线圈中的感应电动势为

$\varepsilon=2Bl_{1}v$

代入数据解得

$\varepsilon=0.31\UV$

（2）$N$ 匝线圈中的总感应电动势为 $N\varepsilon$，由欧姆定律，得

$N\varepsilon=I(R+r)$

代入数字解得

$R+r=12\UO$

Ⅱ

（1）负，阴。

（2）减少。

（3）$N_{A}=\frac{64\Ugmol\times0.21\UA\times3600\Us}{0.25\Ug\times2\times1.6\times10^{-19}\UC}=6.0\times10^{23}\Umoln$
}

\end{enumerate}

\end{document}
